% Unit 058 terjemahan Bahasa Indonesia dari chapter4.tex baris 1681--1840.
% Sumber: Methods of Algebra, Volume 2, commit
% 9a5803ff77dd3257484cb177f851a73770a59dd3 (CC BY 4.0).
% stable-unit-id: o014.aljabr2.chapter4.bounded-derived-functors
% Cakupan: Bagian 4.8 lengkap; berhenti sebelum Bagian 4.9 pada baris 1841.

% segment-id: o014.li2.chapter4.bounded-derived-functors.h001
\section{Funktor Turunan Terbatas}\label{sec:bounded-derived-functor}
\hypertarget{o014.aljabr2.chapter4.bounded-derived-functors}{ }

% segment-id: o014.li2.chapter4.bounded-derived-functors.p001
Bagian ini menjelaskan cara mengkaji ${}^+\mathrm{R}F$ dan
${}^-\mathrm{L}F$ melalui resolusi injektif dan projektif. Kasus kategori
turunan tak terbatas ditunda hingga \S\ref{sec:unbounded-derived}.

% segment-id: o014.li2.chapter4.bounded-derived-functors.c001
\begin{konvensi}
	Bagian ini berfokus pada $\star\in\{+,-\}$ dan menggunakan konvensi berikut.
	\begin{itemize}
		\item Untuk funktor turunan kanan selalu diambil $\star=+$ dan digunakan
		singkatan
		$\mathrm{R}F:={}^+\mathrm{R}F:
		\cate{D}^+(\mathcal{A})\to\cate{D}^+(\mathcal{A}')$.
		\item Untuk funktor turunan kiri selalu diambil $\star=-$ dan digunakan
		singkatan
		$\mathrm{L}F:={}^-\mathrm{L}F:
		\cate{D}^-(\mathcal{A})\to\cate{D}^-(\mathcal{A}')$.
	\end{itemize}
	Berdasarkan Lema~\ref{prop:derived-functor-res}, konvensi ini tidak akan
	menimbulkan kerancuan.
\end{konvensi}

% segment-id: o014.li2.chapter4.bounded-derived-functors.dp001
\begin{definisi-proposisi}\label{prop:F-injective-criterion}
	\index{subkategori tipe I}
	\index{subkategori tipe P}
	Misalkan $\mathcal{A}^\flat$ subkategori penuh aditif dari $\mathcal{A}$.
	Subkategori $\mathcal{A}^\flat$ disebut \emph{bertipe I} relatif terhadap
	$F$ jika syarat-syarat berikut terpenuhi; dalam hal ini
	$\cate{K}^+(\mathcal{A}^\flat)$ merupakan subkategori $F$-injektif dari
	$\cate{K}^+(\mathcal{A})$:
	\begin{enumerate}[({I}1)]
		\item untuk setiap $X\in\Obj(\mathcal{A})$ terdapat monomorfisme
		$X\hookrightarrow I$ dengan $I\in\Obj(\mathcal{A}^\flat)$;
		\item jika $0\to X'\to X\to X''\to0$ merupakan barisan eksak pendek
		dalam $\mathcal{A}$ dan $X',X\in\Obj(\mathcal{A}^\flat)$, maka
		$X''\in\Obj(\mathcal{A}^\flat)$ dan
		$0\to FX'\to FX\to FX''\to0$ eksak.
	\end{enumerate}
	Secara dual, $\mathcal{A}^\flat$ disebut \emph{bertipe P} relatif terhadap
	$F$ jika syarat-syarat berikut terpenuhi; dalam hal ini
	$\cate{K}^-(\mathcal{A}^\flat)$ merupakan subkategori $F$-projektif dari
	$\cate{K}^-(\mathcal{A})$:
	\begin{enumerate}[({P}1)]
		\item untuk setiap $X\in\Obj(\mathcal{A})$ terdapat epimorfisme
		$P\twoheadrightarrow X$ dengan $P\in\Obj(\mathcal{A}^\flat)$;
		\item jika $0\to X'\to X\to X''\to0$ merupakan barisan eksak pendek
		dalam $\mathcal{A}$ dan $X,X''\in\Obj(\mathcal{A}^\flat)$, maka
		$X'\in\Obj(\mathcal{A}^\flat)$ dan
		$0\to FX'\to FX\to FX''\to0$ eksak.
	\end{enumerate}
\end{definisi-proposisi}
% segment-id: o014.li2.chapter4.bounded-derived-functors.pf001
\begin{bukti}
	Cukup dibahas versi tipe I. Pertama,
	$\cate{K}^+(\mathcal{A}^\flat)$ merupakan subkategori bertriangulasi dari
	$\cate{K}^+(\mathcal{A})$ (Korolari~\ref{prop:subcat-K-triangulated}).
	Teorema~\ref{prop:resolution-existence} menunjukkan bahwa untuk setiap
	$X\in\Obj(\cate{K}^+(\mathcal{A}))$ terdapat kuasi-isomorfisme $X\to I$
	dengan $I\in\cate{K}^+(\mathcal{A}^\flat)$. Ini adalah kondisi resolusi
	bagi subkategori $F$-injektif. Tinggal dibuktikan bahwa jika
	$X\in\cate{K}^+(\mathcal{A}^\flat)$ asiklik, maka $FX$ juga asiklik.

	Untuk setiap $n\in\Z$ terdapat barisan eksak pendek
	$0\to\Ker(d^n)\to X^n\xrightarrow{d^n}\Ker(d^{n+1})\to0$. Jika $n\ll0$,
	$\Ker(d^n)=X^n=0$. Dengan menerapkan (I2) secara rekursif untuk semua $n$,
	diperoleh $\Ker(d^n)\in\Obj(\mathcal{A}^\flat)$ serta barisan eksak pendek
	dalam $\mathcal{A}'$
	\[ 0 \to F\Ker(d^n) \to FX^n \xrightarrow{Fd^n} F\Ker(d^{n+1}) \to 0. \]
	Penyambungan barisan-barisan eksak pendek ini menunjukkan bahwa $FX$
	asiklik.
\end{bukti}

% segment-id: o014.li2.chapter4.bounded-derived-functors.p002
Istilah tipe I dan tipe P hanya merupakan sebutan praktis. Terdapat
karakterisasi lengkap bagi subkategori $\mathcal{A}^\flat$ yang memberikan
subkategori $F$-injektif atau $F$-projektif; lihat
\cite[Proposisi 13.3.5]{KS06}.

% segment-id: o014.li2.chapter4.bounded-derived-functors.p003
Hasil berikut menunjukkan bahwa funktor turunan dapat dihitung dengan
mengambil resolusi di dalam subkategori penuh bertipe I atau P.

% segment-id: o014.li2.chapter4.bounded-derived-functors.t001
\begin{teorema}\label{prop:derived-via-resolution}
	Misalkan terdapat subkategori penuh aditif $\mathcal{A}^\flat$ dari
	$\mathcal{A}$ yang bertipe I (atau tipe P) relatif terhadap funktor aditif
	$F:\mathcal{A}\to\mathcal{A}'$. Maka funktor turunan $\mathrm{R}F$ (atau
	$\mathrm{L}F$) ada. Lebih tepatnya:
	\begin{enumerate}[(i)]
		\item untuk setiap $X\in\Obj(\cate{K}^+(\mathcal{A}))$ (atau
		$X\in\Obj(\cate{K}^-(\mathcal{A}))$), terdapat kuasi-isomorfisme
		$X\to I$ (atau $P\to X$) dengan
		$I\in\Obj(\cate{K}^+(\mathcal{A}^\flat))$ (atau
		$P\in\Obj(\cate{K}^-(\mathcal{A}^\flat))$);
		\item bagi kuasi-isomorfisme pada (i), morfisme dalam
		\eqref{eqn:derived-can-morphism} memberikan
		\[ \mathrm{R}F(QX) \rightiso \mathrm{R}F(QI) \leftiso Q' \cate{K}^+ F(I), \quad \mathrm{L}F(QX) \leftiso \mathrm{L}F(QP) \rightiso Q' \cate{K}^- F(P); \]
		\item $\mathrm{R}F$ membatasi diri pada
		$\cate{D}^{\geq0}(\mathcal{A})\to\cate{D}^{\geq0}(\mathcal{A}')$
		(atau $\mathrm{L}F$ membatasi diri pada
		$\cate{D}^{\leq0}(\mathcal{A})\to\cate{D}^{\leq0}(\mathcal{A}')$);
		\item misalkan $F$ eksak kiri (atau eksak kanan). Jika
		$X\in\Obj(\mathcal{A})$, morfisme dalam
		\eqref{eqn:derived-can-morphism} memberikan isomorfisme
		$FX\rightiso\mathrm{R}^0F(QX)$ (atau
		$FX\leftiso\mathrm{L}_0F(QX)$);
		\item jika $X\in\Obj(\mathcal{A}^\flat)$, maka untuk $n>0$ berlaku
		$\mathrm{R}^nF(X)=0$ (atau $\mathrm{L}_nF(X)=0$).
	\end{enumerate}
\end{teorema}
% segment-id: o014.li2.chapter4.bounded-derived-functors.pf002
\begin{bukti}
	Cukup dibahas kasus $\mathrm{R}F$. Karena
	$\cate{K}^+(\mathcal{A}^\flat)$ bersifat $F$-injektif, kuasi-isomorfisme
	$X\to I$ selalu ada. Keberadaan $\mathrm{R}F$ dan isomorfisme
	$\mathrm{R}F(QX)\simeq Q'\cate{K}^+F(I)$ hanyalah pernyataan ulang
	Proposisi~\ref{prop:localization-injective}. Ini membuktikan (i) dan (ii).

	Untuk (iii), jika $X\in\Obj(\cate{D}^{\geq0}(\mathcal{A}))$, kita boleh
	menganggap bahwa $X$ berasal dari objek $\cate{C}^{\geq0}(\mathcal{A})$.
	Teorema~\ref{prop:resolution-existence} memungkinkan pemilihan
	kuasi-isomorfisme $X\to I$ dengan
	$I\in\Obj(\cate{C}^{\geq0}(\mathcal{A}^\flat))$. Jadi
	\[ \mathrm{R}F(QX) \simeq Q' \cate{K}^+ F(I) \in \Obj\left( \cate{D}^{\geq 0}(\mathcal{A}')\right). \]
	Jika juga $X\in\Obj(\mathcal{A}^\flat)$, kuasi-isomorfisme $X\to I$ dapat
	diambil sebagai $\identity_X$. Dalam hal ini
	$\mathrm{R}F(QX)=Q'\cate{K}^+F(X)=Q'FX$ terkonsentrasi pada derajat nol.
	Ini sekaligus membuktikan (v).

	Terakhir, tinjau (iv). Misalkan $F$ eksak kiri. Untuk
	$X\in\Obj(\mathcal{A})$, pilih barisan eksak
	$0\to X\to I^0\to I^1\to\cdots$ dengan setiap
	$I^n\in\Obj(\mathcal{A}^\flat)$. Ambil $I=[I^0\to I^1\to\cdots]$.
	Dari (ii) dan keeksakan kiri diperoleh
	$FX\rightiso\Ker[Fd_I^0:F(I^0)\to F(I^1)]
	\simeq\mathrm{R}^0F(X)$, sebagaimana dikehendaki.
\end{bukti}

% segment-id: o014.li2.chapter4.bounded-derived-functors.e001
\begin{contoh}[Objek injektif dan projektif]\label{eg:injective-gives-F-injectives}
	Misalkan $\mathcal{A}$ mempunyai cukup banyak objek injektif, dan ambil
	$\mathcal{A}^\flat$ sebagai subkategori penuh aditif yang terdiri atas objek
	injektif. Subkategori ini bertipe I relatif terhadap setiap $F$: syarat (I1)
	terpenuhi menurut asumsi, sedangkan untuk (I2),
	Lema~\ref{prop:inj-proj-split} menunjukkan bahwa barisan eksak pendek dalam
	syarat tersebut terbelah. Jadi $X'\oplus X''$ injektif, dan
	Lema~\ref{prop:prod-injective-projective} menyiratkan bahwa $X''$ juga
	injektif.

	Secara dual, jika $\mathcal{A}$ mempunyai cukup banyak objek projektif,
	subkategori penuh aditif yang terdiri atas objek-objek tersebut bertipe P
	relatif terhadap setiap $F$.
\end{contoh}

% segment-id: o014.li2.chapter4.bounded-derived-functors.p004
Contoh~\ref{eg:injective-gives-F-injectives} menghubungkan pembahasan ini
dengan definisi klasik funktor turunan dalam \S\ref{sec:derived-primer}.
Misalkan $\mathcal{A}$ mempunyai cukup banyak objek injektif. Tinjau pembatasan
funktor turunan kanan $\mathrm{R}^nF|_{\mathcal{A}}$. Karena barisan eksak
pendek dalam $\mathcal{A}$ secara kanonik dapat diperluas menjadi segitiga
terbedakan (Proposisi~\ref{prop:ses-vs-triangle}), barisan eksak panjang segera
menjadikan $(\mathrm{R}^nF|_{\mathcal{A}})_{n\geq0}$ suatu funktor-$\delta$
kohomologis dalam pengertian Definisi~\ref{def:delta-functor}. Secara dual,
funktor turunan kiri memberikan funktor-$\delta$ homologis
$(\mathrm{L}_nF|_{\mathcal{A}})_{n\geq0}$.

% segment-id: o014.li2.chapter4.bounded-derived-functors.pr001
\begin{proposisi}
	Misalkan $F$ eksak kiri (atau eksak kanan), dan $\mathcal{A}$ mempunyai cukup
	banyak objek injektif (atau projektif). Maka
	$\mathrm{R}^nF:\cate{C}^+(\mathcal{A})\to\mathcal{A}'$ (atau
	$\mathrm{L}_nF:\cate{C}^-(\mathcal{A})\to\mathcal{A}'$) sama dengan versi
yang didefinisikan sebelumnya dalam Definisi~\ref{def:derived-primer}.

	Lebih umum, jika terdapat subkategori $\mathcal{A}^\flat$ bertipe I (atau
	tipe P), sebagaimana dalam
	Definisi--Proposisi~\ref{prop:F-injective-criterion}, maka
	$(\mathrm{R}^nF|_{\mathcal{A}})_{n\geq0}$ (atau
	$(\mathrm{L}_nF|_{\mathcal{A}})_{n\geq0}$) merupakan funktor-$\delta$
	universal dalam pengertian Definisi~\ref{def:universal-delta-functor}.
\end{proposisi}
% segment-id: o014.li2.chapter4.bounded-derived-functors.pf003
\begin{bukti}
	Cukup dibahas kasus eksak kiri. Untuk
	$X\in\Obj(\cate{C}^+(\mathcal{A}))$, pilih resolusi injektif $X\to I$.
	Teorema~\ref{prop:derived-via-resolution} (ii) memberikan isomorfisme
	kanonik
	$\mathrm{R}^nF(QX)\simeq\Hm^n\cate{K}^+F(I)$ untuk setiap $n\in\Z$.
	Ruas kanan tepat merupakan funktor turunan versi
	Definisi~\ref{def:derived-primer}.

	Lebih umum, misalkan $\mathcal{A}^\flat$ bertipe I relatif terhadap $F$.
	Untuk membuktikan universalitas
	$(\mathrm{R}^nF|_{\mathcal{A}})_{n\geq0}$, berdasarkan
	Proposisi~\ref{prop:erasable-univ-delta} cukup dibuktikan bahwa
	$\mathrm{R}^nF|_{\mathcal{A}}$ dapat dihapus untuk $n>0$. Ini langsung
	mengikuti syarat (I1) dalam
	Definisi--Proposisi~\ref{prop:F-injective-criterion} dan
	Teorema~\ref{prop:derived-via-resolution} (v).
\end{bukti}

% segment-id: o014.li2.chapter4.bounded-derived-functors.p005
Untuk menurunkan komposisi funktor, sering diperlukan kelas objek yang lebih
besar sebagai resolusi. Kita memerlukan konsep berikut.

% segment-id: o014.li2.chapter4.bounded-derived-functors.d001
\begin{definisi}\label{def:F-acyclic-object}
	\index{objek asiklik-$F$}
	Misalkan $\mathrm{R}F$ (atau $\mathrm{L}F$) ada. Suatu objek
	$X\in\Obj(\mathcal{A})$ disebut \emph{$F$-asiklik} jika
	$\mathrm{R}F(X)$ (atau $\mathrm{L}F(X)$) merupakan objek dari
	$\cate{D}^{\geq0}(\mathcal{A}')\cap\cate{D}^{\leq0}(\mathcal{A}')
	\simeq\mathcal{A}'$.
\end{definisi}

% segment-id: o014.li2.chapter4.bounded-derived-functors.p006
Hasil berikut menunjukkan bahwa funktor turunan dapat dihitung melalui
resolusi $F$-asiklik; bandingkan dengan Korolari~\ref{prop:acyclic-resolution}.

% segment-id: o014.li2.chapter4.bounded-derived-functors.co001
\begin{korolari}\label{prop:derived-via-F-acyclic}
	Misalkan subkategori penuh aditif $\mathcal{A}^\flat$ dari $\mathcal{A}$
	bertipe I (atau tipe P) relatif terhadap $F$. Definisikan
	$\mathcal{A}^\natural$ sebagai subkategori penuh aditif yang terdiri atas
	semua objek $F$-asiklik dalam $\mathcal{A}$. Maka:
	\begin{enumerate}[(i)]
		\item $\mathcal{A}^\natural\supset\mathcal{A}^\flat$;
		\item $\mathcal{A}^\natural$ juga bertipe I (atau tipe P);
		\item untuk setiap kuasi-isomorfisme $X\to I$ (atau $P\to X$), dengan
		$I\in\Obj(\cate{K}^+(\mathcal{A}^\natural))$ (atau
		$P\in\Obj(\cate{K}^-(\mathcal{A}^\natural))$, morfisme dalam
		\eqref{eqn:derived-can-morphism} memberikan
		\[ \mathrm{R}F(QX) \rightiso \mathrm{R}F(QI) \leftiso Q' \cate{K}^+ F(I), \quad \mathrm{L}F(QX) \leftiso \mathrm{L}F(QP) \rightiso Q' \cate{K}^- F(P). \]
	\end{enumerate}
\end{korolari}
% segment-id: o014.li2.chapter4.bounded-derived-functors.pf004
\begin{bukti}
	Teorema~\ref{prop:derived-via-resolution} (v) menyiratkan (i), sehingga
	syarat (I1) atau (P1) juga berlaku bagi $\mathcal{A}^\natural$. Untuk (I2),
	misalkan $0\to X'\to X\to X''\to0$ barisan eksak pendek dalam
	$\mathcal{A}$ dan $X',X$ bersifat $F$-asiklik. Barisan eksak panjang dari
	Teorema~\ref{prop:derived-long-exact} memberikan barisan eksak
	\[ \underbracket{\mathrm{R}^n F(X)}_{= 0} \to \mathrm{R}^n F(X'') \to \underbracket{\mathrm{R}^{n+1}F(X')}_{= 0}, \quad n \in \Z_{\geq 1} , \]
	sehingga $X''$ bersifat $F$-asiklik. Di sisi lain, bagian $n=0$ dari barisan
	eksak panjang bersama $\mathrm{R}^1F(X')=0$ menunjukkan bahwa
	$0\to FX'\to FX\to FX''\to0$ eksak. Pembuktian (P2) bersifat dual. Ini
	membuktikan (ii).

	Terakhir, (iii) diperoleh dengan menerapkan
	Teorema~\ref{prop:derived-via-resolution} (ii) pada $\mathcal{A}^\natural$.
\end{bukti}

% segment-id: o014.li2.chapter4.bounded-derived-functors.p007
Sekarang kita membahas penurunan komposisi funktor. Alat dasarnya adalah
Teorema~\ref{prop:localization-triangulated-composite}.

% segment-id: o014.li2.chapter4.bounded-derived-functors.t002
\begin{teorema}[Menurunkan komposisi funktor]\label{prop:derived-composite}
	Tinjau funktor aditif antarkategori abelian
	\[ \mathcal{A} \xrightarrow{F} \mathcal{A}' \xrightarrow{F'} \mathcal{A}''. \]
	Misalkan terdapat subkategori penuh $\mathcal{A}^\flat$ dari $\mathcal{A}$
	dan $(\mathcal{A}')^\flat$ dari $\mathcal{A}'$, yang masing-masing bertipe I
	relatif terhadap $F$ dan $F'$ dalam pengertian
	Definisi--Proposisi~\ref{prop:F-injective-criterion}. Jika $F$ memetakan
	objek-objek $\mathcal{A}^\flat$ ke objek $F'$-asiklik dalam $\mathcal{A}'$,
	maka morfisme kanonik
	$\mathrm{R}(F'F)\to(\mathrm{R}F')(\mathrm{R}F)$ merupakan isomorfisme.

	Terdapat versi yang bersesuaian bagi funktor turunan kiri $\mathrm{L}F$,
	$\mathrm{L}F'$, dan morfisme kanonik
	$(\mathrm{L}F')(\mathrm{L}F)\to\mathrm{L}(F'F)$, dengan menggunakan
	subkategori bertipe P.
\end{teorema}
% segment-id: o014.li2.chapter4.bounded-derived-functors.pf005
\begin{bukti}
	Cukup dibahas kasus funktor turunan kanan. Misalkan
	$(\mathcal{A}')^\natural$ subkategori penuh aditif yang terdiri atas objek
	$F'$-asiklik dalam $\mathcal{A}'$.
	Korolari~\ref{prop:derived-via-F-acyclic} bersama
	Definisi--Proposisi~\ref{prop:F-injective-criterion} menunjukkan bahwa
	$\cate{K}^+(\mathcal{A}^\flat)$ dan
	$\cate{K}^+((\mathcal{A}')^\natural)$ berturut-turut merupakan subkategori
	$F$-injektif dan $F'$-injektif dari $\cate{K}^+(\mathcal{A})$ dan
	$\cate{K}^+(\mathcal{A}')$. Terapkan
	Teorema~\ref{prop:localization-triangulated-composite}.
\end{bukti}

% segment-id: o014.li2.chapter4.bounded-derived-functors.p008
Dengan menggunakan amplitudo dari Definisi~\ref{def:functor-amplitude}, kita
dapat mendefinisikan dimensi funktor turunan kanan dan kiri $F$.

% segment-id: o014.li2.chapter4.bounded-derived-functors.dp002
\begin{definisi-proposisi}\label{def:dim-functor}
	\index{dimensi}
	Misalkan $F:\mathcal{A}\to\mathcal{A}'$ eksak kiri (atau eksak kanan), dan
	funktor turunan $\mathrm{R}F$ (atau $\mathrm{L}F$) ada serta ditentukan
seperti dalam Teorema~\ref{prop:derived-via-resolution}. Terdapat
$d\in\Z_{\geq0}\sqcup\{+\infty\}$ sehingga amplitudo funktor turunan tersebut
termuat dalam $[0,d]$ (atau $[-d,0]$). Infimum semua $d$ demikian disebut
\emph{dimensi} $\mathrm{R}F$ (atau $\mathrm{L}F$).
\end{definisi-proposisi}
% segment-id: o014.li2.chapter4.bounded-derived-functors.pf006
\begin{bukti}
	Cukup dibahas kasus eksak kiri. Teorema~\ref{prop:derived-via-resolution}
	menyiratkan bahwa $\mathrm{R}F$ membatasi diri pada
	$\mathcal{A}\to\cate{D}^{\geq0}(\mathcal{A}')$, sehingga amplitudonya
	termuat dalam suatu interval $[0,d]$.
\end{bukti}

% segment-id: o014.li2.chapter4.bounded-derived-functors.p009
	Sebagai akibat, jika $\mathrm{R}F$ (atau $\mathrm{L}F$) ada dan berdimensi
	hingga, maka funktor tersebut membatasi diri pada
$\cate{D}^{\bdd}(\mathcal{A})\to\cate{D}^{\bdd}(\mathcal{A}')$.

% segment-id: o014.li2.chapter4.bounded-derived-functors.e002
\begin{contoh}[Dimensi $\Tor$]\label{eg:Tor-dimension}
	\index{dimensi!$\operatorname{Tor}$}
	Misalkan $R$ gelanggang dan $X$ modul kanan $R$. Tinjau funktor eksak kanan
	dari $R\dcate{Mod}$ ke $\cate{Ab}$ yang memetakan
	$Y\mapsto X\otimes Y:=X\dotimes{R}Y$. Karena $R\dcate{Mod}$ mempunyai cukup
	banyak objek projektif, $\mathrm{L}(X\otimes\cdot)$ ada, dan
	$\mathrm{L}_n(X\otimes\cdot)$ tepat merupakan
	$\Tor^R_n(X,\cdot)$ dari Definisi--Proposisi~\ref{def:Tor-classical}.
	Dimensi $\mathrm{L}(X\otimes\cdot)$ disebut \emph{dimensi $\Tor$} dari
	$X$. Definisi serupa berlaku bagi modul kiri $R$, dan
	Proposisi~\ref{prop:Tor-symmetry} menyiratkan bahwa keduanya sama jika $R$
	komutatif. Funktor hasil kali tensor turunan dibahas lebih lanjut dalam
	\S\ref{sec:otimesL}.
\end{contoh}

% segment-id: o014.li2.chapter4.bounded-derived-functors.p010
Sekarang tetapkan kategori abelian $\mathcal{A}_1$, $\mathcal{A}_2$,
$\mathcal{B}$ dan bifunktor aditif
\[ F: \mathcal{A}_1 \times \mathcal{A}_2 \to \mathcal{B}. \]
Kita tidak mensyaratkan $\mathcal{B}$ mempunyai koproduk atau produk terhitung,
sebab $\cate{C}^2F$ membatasi menjadi
\begin{gather*}
	\cate{C}^2 F: \cate{C}^{\pm}(\mathcal{A}_1) \times \cate{C}^{\pm} (\mathcal{A}_2) \to \cate{C}^2_f(\mathcal{B}).
\end{gather*}
Pada $\cate{C}^2_f(\mathcal{B})$, funktor kompleks total
$\tot_{\oplus}=\tot_{\Pi}$ selalu terdefinisi; tuliskan funktor ini sebagai
$\tot$. Demikian pula, tuliskan
$\cate{C}_{\oplus}F=\cate{C}_{\Pi}F$ sebagai $\cate{C}F$, dan
$\cate{K}_{\oplus}F=\cate{K}_{\Pi}F$ sebagai $\cate{K}F$.

% segment-id: o014.li2.chapter4.bounded-derived-functors.p011
Untuk mengkaji bifunktor turunan dalam Definisi~\ref{def:derived-bifunctor},
kita meninjau subkategori bertriangulasi $\mathcal{I}_i$ (atau
$\mathcal{P}_i$) dari $\cate{K}^+(\mathcal{A}_i)$ (atau
$\cate{K}^-(\mathcal{A}_i)$). Seperti biasa, tanda ditentukan oleh apakah
yang dibahas bifunktor turunan kanan atau kiri, dan $i=1,2$. Bagaimana pasangan
$(\mathcal{I}_1,\mathcal{I}_2)$ atau $(\mathcal{P}_1,\mathcal{P}_2)$ dipilih?
Definisi--Proposisi~\ref{prop:F-injective-criterion} dan teknik dalam
\S\ref{sec:double-cplx-coh} memberikan jawaban berikut.

% segment-id: o014.li2.chapter4.bounded-derived-functors.l001
\begin{lema}\label{prop:RF-bifunctor}
	Misalkan $\mathcal{A}_i^\flat\subset\mathcal{A}_i$ subkategori penuh aditif,
	$i=1,2$. Jika syarat-syarat berikut berlaku sekaligus, maka pasangan
	$\bigl(\cate{K}^+(\mathcal{A}_1^\flat),
	\cate{K}^+(\mathcal{A}_2^\flat)\bigr)$ (atau
	$\bigl(\cate{K}^-(\mathcal{A}_1^\flat),
	\cate{K}^-(\mathcal{A}_2^\flat)\bigr)$) bersifat $F$-injektif (atau
	$F$-projektif) dalam pengertian
	Definisi~\ref{def:F-injective-bifunctor}:
	\begin{itemize}
		\item jika $X_1\in\Obj(\mathcal{A}_1^\flat)$ tetap, maka
		$\mathcal{A}_2^\flat$ bertipe I (atau tipe P) relatif terhadap
		$F(X_1,\cdot)$;
		\item jika $X_2\in\Obj(\mathcal{A}_2^\flat)$ tetap, maka
		$\mathcal{A}_1^\flat$ bertipe I (atau tipe P) relatif terhadap
		$F(\cdot,X_2)$.
	\end{itemize}
	Sebagai kasus khusus, jika $\mathcal{A}_1$ dan $\mathcal{A}_2$ masing-masing
	mempunyai cukup banyak objek injektif (atau projektif), maka
	$\mathrm{R}F$ (atau $\mathrm{L}F$) ada.
\end{lema}
% segment-id: o014.li2.chapter4.bounded-derived-functors.pf007
\begin{bukti}
	Ambil $\mathcal{I}_i:=\cate{K}^+(\mathcal{A}_i^\flat)$, $i=1,2$. Satu-satunya
	hal yang perlu dibuktikan adalah bahwa untuk setiap
	$X_1\in\Obj(\cate{K}^+(\mathcal{A}_1^\flat))$, funktor
	$(\cate{K}F)(X_1,\cdot):\cate{K}^+(\mathcal{A}_2^\flat)
	\to\cate{K}^+(\mathcal{B})$ memetakan kompleks asiklik ke kompleks
	asiklik; hal yang sama berlaku setelah menukar indeks $1$ dan $2$.
	Keberadaan $\mathrm{R}F$ (atau $\mathrm{L}F$) kemudian langsung mengikuti
	dari Proposisi~\ref{prop:localization-injective-bifunctor}.

	Ambil $Y\in\Obj(\cate{C}^+(\mathcal{A}_2^\flat))$. Maka
	$(\cate{K}^2F)(X_1,Y)$ berada dalam
	$\Obj(\cate{C}^2_f(\mathcal{B}))$. Menurut syarat (I2), setiap kolom
	$F(X_1^p,Y^\bullet)$ asiklik. Karena itu,
	Korolari~\ref{prop:acyclic-assembly} menyiratkan bahwa
	$\cate{K}F(X_1,Y):=\tot(\cate{K}^2F(X_1,Y))$ eksak, sebagaimana
	dikehendaki.
\end{bukti}

% segment-id: o014.li2.chapter4.bounded-derived-functors.p012
Jika terdapat subkategori $F$-injektif (atau $F$-projektif), funktor
$\mathrm{R}F$ (atau $\mathrm{L}F$) dapat ditentukan menurut
\eqref{eqn:derived-bifunctor-1} (atau \eqref{eqn:derived-bifunctor-2}).
